% ============================================================
%  CHAPTER II
%  Preliminary Remarks on Visual Perception
%  Source pp. 42–51
% ============================================================
\chapter{Preliminary Remarks on Visual Perception}
\markboth{Chapter II. Preliminary Remarks on Visual Perception}{Chapter II}

\srcpage{42}

Man lives in three-dimensional space, which is populated by numerous objects and
various living creatures~--- some useful, some dangerous, some indifferent.  Visual
perception of the surrounding space warns of these objects and creatures long before
direct contact with them, which is of fundamental importance to the behaviour of
animals and human beings alike.  Naturally, the quality and reliability of this
information are of exceptional significance in the biological struggle for existence.

How faithfully does the visual impression of the surrounding world correspond to
reality?  The theory of reflection provides the answer: the psyche is a reflection of
objectively existing reality, and perception is only the first, directly sensory form of
that reflection.  This first form of reflection is not passive or purely mirror-like;
perception is the result of active human activity, and even visual perception contains,
as an essential constituent, the active engagement of the perceiving subject.

If one schematizes the process of visual perception, the simplest model describes it
as two-staged.  The first stage is the formation of an image of the external space on the
retina of the eye; the second is the reconstruction, on this basis, of the appearance of
the external space in human consciousness.  That these two stages are qualitatively
distinct is apparent even from the geometric fact that the image on the retina is
two-dimensional (approximately flat, in first approximation, like a picture), while the
representation of the external space formed on that basis is three-dimensional.\footnote{%
  This chapter does not attempt to survey current views on the visual perception of
  space by human beings.  Only those aspects of the problem that are indispensable for
  understanding this book will be briefly considered below.  A sufficiently complete and
  accessible account of the current state of the question may be found in the books of
  R.\,L.~Gregory (\emph{Eye and Brain}. Moscow, 1970; \emph{The Intelligent Eye}.
  Moscow, 1972).}

The formation of a two-dimensional image on the retina is a process sufficiently well
understood, and in its essentials analogous to the formation of an image on
photographic film by a camera lens.  There are, to be sure, psychophysiologically
significant differences: the eye is in continuous motion, ``scanning'' the external space
with a comparatively narrow beam of vision, and studies of this subconscious movement
demonstrate its exceptional importance in the process of visual perception.  However,
since in the questions considered below the most important issue is the geometric
properties of the image formed on the retina, the simplest model will be used
throughout: on the retina, within a sufficiently large angle of vision, a two-dimensional
image of the external world is obtained by projecting it through the lens, which plays a
role analogous to that of a camera objective.  Such an image on the retina of the eye is
called the \emph{retinal image} (from \textit{retina}~--- the screen of the eye).\footnote{%
  It is worth reemphasizing that the concept of ``retinal image'' as used here and
  throughout is a sufficiently conventional one: it takes no account of the actual
  curvature of the retinal surface or of the mobility of the eye.  Moreover, the retina
  itself is not homogeneous~--- the central (foveal) part of the retina gives the best
  perception of detail and colour, while its peripheral part yields low acuity.
  Consequently, the concept of ``retinal image'' as used in this book is a certain
  idealization, accounting for the aggregate effect of the first stage of visual
  perception.}

\srcpage{43}

The retinal image gives a strongly distorted picture of the external world.  These
distortions arise above all from the fact that nearby objects appear enormous while more
distant ones appear small.

The retinal image in itself contains no information about the distances to the objects
whose images have formed on the retina.  The necessary supplementation of the retinal
image is therefore achieved by drawing on other information.  This process of
introducing supplementary information constitutes the second stage of visual
perception.  The required information is drawn from the ``memory reserves'' of the
human brain.  Without undue oversimplification, one may assert that human
consciousness stores in its memory the experience of previous life from the first days of
infancy, and even genetically transmitted experience of preceding generations.  Using
this experience, human consciousness reconstructs from the retinal image a true picture
of the external world with a reasonable degree of accuracy.  For what follows it is
important to note that the image of the external world thus created by the brain is not an
exact copy of the real external world.  Such exactness would be biologically pointless
and even harmful~--- the perception of nearby regions of space, which may conceal
danger or serve for obtaining food, must be accurate, while more distant regions can and
even should reasonably be perceived with less detail and precision (so as not to overload
the brain with information inessential to the determination of behaviour).

These ``memory reserves'' are nothing other than the experience, acquired through
active human activity, that permits the surrounding objective world to be perceived
correctly to a reasonable approximation~--- an adequate, if approximate, reflection of
it in human consciousness.  The experience of knowing the external world referred to
here is not the sum or result of accumulated visual information; it encompasses all
of human practical activity, which allows one's representations of the external world
to be corrected, correlated with visual perception, and ultimately brought to constitute
an adequate visual image of the surrounding objective world with the requisite degree
of detail and accuracy.

\srcpage{44}

The image of the external world reconstructed by the brain from the retinal image
acquires the form of a three-dimensional space, but, as already noted, one differing
from the real three-dimensional space.  Since an artist whose goal is the making of a
picture~--- as opposed to a technical drawing~--- must be able to depict precisely this
visually perceived (perceptual) space, its properties merit closer examination.

Confining attention to geometry, the problem reduces to the question of how the
two-dimensional retinal image is transformed into three-dimensional perceptual space,
or more precisely, what geometric laws govern this transformation.

First of all (not in time, of course, but logically), all the images recorded on the retina
must be mentally distributed at their corresponding distances from the person
contemplating them.  This distribution along the ``near--far'' axis occurs subconsciously,
on the basis of prior experience, the person making subconscious use of a number of
\emph{depth cues} of the visible space.  These cues may be divided, somewhat
conventionally, into two groups~--- monocular and binocular.  The first group comprises
those depth cues that are effective even when one looks with one eye.

The monocular depth cues include the following:
(1)~\emph{occlusion}~--- nearby objects can block more distant ones from view;
(2)~\emph{reduction in size} of objects with increasing distance from the viewer~---
particularly effective for objects whose true size is known (trees, etc.);
(3)~\emph{aerial perspective}~--- distant objects are seen less sharply and change
colour (appearing, as it were, through a bluish haze);
(4)~\emph{position in the visual field}~--- distant objects are seen shifted upward in the
field of vision, that is, closer to the horizon;
(5)~\emph{shadow and lighting}~--- large objects (round fortification towers, large
buildings, etc.) illuminated by the sun are shaded differently in their different parts,
and the observed pattern of shadows can give an impression of the nearer and more
distant portions of these objects; this cue is particularly effective for large curvilinear
and ridged surfaces.

Not all monocular depth cues have been listed here, only the principal ones, namely
those that will be needed below.

The binocular depth cues include:
(1)~\emph{convergence}~--- the rotation of the optical axes of the eyes toward the object
being examined; normally (when looking into the distance) the optical axes of the eyes
are parallel, but if a sufficiently close object is examined, the corresponding muscles
converge the optical axes of the eyes so that they intersect in the region being
scrutinized, and the corresponding muscular effort signals the nearness or remoteness of
the object;
(2)~\emph{disparity}~--- even without actively looking at some object, since the eyes are
spaced apart, the image received by the left and right eye differs, and the more so the
closer the objects are to the viewer.

\srcpage{45}

Acting together, the depth cues provide a sufficiently complete impression of the
distances to various objects; moreover, as objects approach the viewer, the number and
``accuracy'' of the cues increase in full accord with the biological needs of the human
being.  For large distances, the required information about distances to objects is
provided by the first four monocular cues, while in the immediate vicinity of a person
the binocular cues and all the monocular cues are active~--- with the exception only of
aerial perspective.

Careful consideration of all the depth cues listed shows that they cannot be obtained
directly from the retinal image, but arise from prior experience.

Even such an apparently unconditional cue as the blocking of a distant object by a
nearby one is not absolute: it presupposes knowledge of the true form of the partially
blocked object, for otherwise it is impossible to determine whether its visible portion is
the complete (hence unblocked) form.  Thus the human brain, using accumulated prior
experience, provides the fundamental possibility of constructing a three-dimensional
perceptual space from the two-dimensional retinal image.

In addition to the subconscious analysis of visual information, the formation of a
spatial image in human consciousness is assisted by information (likewise drawn from
memory reserves) that has nothing to do with vision.  Living in a familiar world, a person
not only contemplates it but continually interacts with it in practical activity: moving
about, picking up objects, and so forth.  He therefore knows excellently the dimensions
and shape of his room, his writing desk, the objects lying on it, and the geometric
properties of his surroundings in general.  That this knowledge can be entirely
unconnected with vision (though greatly helpful to it) is shown by the practical activity
of the blind: a blind person unerringly finds the objects he needs in his room, quickly
and without assistance moves about his flat, and so on.  This knowledge of space, its
dimensions and form, may be called a \emph{tactile-kinesthetic representation} of it,
allowing in particular the distances of objects from a person to be known.

Although in somewhat different form and possibly somewhat diminished, this
knowledge is also present in the sighted.  This is easy to verify by examining nearby
objects first with both eyes and then closing one eye.  One might expect that
``disconnecting'' the binocular depth cues, which are so important at small distances,
would cause a marked change in visual perception.  Yet nothing of the sort occurs: the
person continues to see objects essentially the same, and~--- importantly~--- as
three-dimensional.  Tactile-kinesthetic representations take the place of the binocular
depth cues in this situation.  Thus binocular depth cues and tactile-kinesthetic
representations act, as it were, in a kind of unity, contributing to accurate assessment
of distances to objects.

\srcpage{46}

In cases where both these sources of information were simultaneously ``disconnected''
artificially~--- i.e., when a person, observing with one eye from a suitably chosen
vantage point, was presented with spatial constructions of forms unknown and
unfamiliar to him~--- it proved possible to induce in the person visual illusions that
completely distorted the actual spatial relationships.  This is demonstrated in particular
by the well-known experiments with Ames's ``distorted room.''~\footnote{%
  See, for example: R.\,L.~Gregory. \emph{Eye and Brain}. pp.\,194--198; idem.
  \emph{The Intelligent Eye}. pp.\,27--29.}

Thus the depth cues available to a person for ``placing'' the objects recorded in the
retinal image along the near--far axis reduce to monocular and binocular cues,
supplemented by tactile-kinesthetic representations.

The effectiveness of tactile-kinesthetic representations, like that of the binocular
depth cues, declines with increasing distance to the objects contemplated, once
distances exceed the somewhat indefinite boundary of the person's ``immediate
surroundings.''  Large distances are not so biologically important, and therefore an
excessive loading of memory with detailed information about large regions of space
would be biologically unreasonable.

Returning to the question of how perceptual space is formed in human consciousness,
let us note an almost obvious circumstance: merely ``placing'' the objects registered by
the retina along some distance scale is insufficient.  The optical processes associated
with the action of the eye produce on the retina a manifestly distorted picture of the
external world~--- nearby objects appear large and distant ones small, even if in
objective external space they are quite identical.  Such distortion can be tolerated for
distant regions of space, but for nearby ones, which are of primary biological
importance, distortions of any kind are highly undesirable, as they can lead to errors in
behaviour.  Therefore one of the tasks of the visual perception system is to rework, in
the course of the second stage of perception, the geometric relationships of the retinal
image~--- a reworking that is stronger the nearer the region of space being
contemplated.  These processes of ``correcting'' the geometry of the retinal image are
now well understood and have received in the psychology of visual perception the name
\emph{constancy mechanisms}.  We shall dwell on two such mechanisms~--- the
\emph{size-constancy mechanism} and the \emph{form-constancy mechanism}. (Other
constancy mechanisms~--- colour constancy, for example~--- will not be discussed, as
they are not geometric in character.)

The size-constancy mechanism is concerned with compensating for the reduction in
the image of an object on the retina as that object recedes.  In the zone of a person's
immediate surroundings (i.e.\ within a radius of a few metres) this compensation is
nearly complete.  Consequently, the visual image of a nearby object arising in one's
consciousness may differ substantially in apparent size from the corresponding retinal
image, yet will be in agreement with the true size of the object contemplated.  This is
easy to verify by looking at the floor of a small room: observing the floorboards or
parquet strips at a distance of three to four metres and near one's feet, one hardly
perceives any visible narrowing, whereas on the retina of the eye the corresponding
narrowing would be approximately twofold (in agreement with the ratio of the named
distance to the height of the person).

\srcpage{47}

Thus, although the size of an object on the retina diminishes proportionally to distance
as the object recedes, the perceived size remains nearly unchanged, constant~--- hence
the name: \emph{size-constancy mechanism}.  This fact is well known to portrait
painters.  When creating a group portrait, the artist paints the heads of all the sitters
at approximately the same size, even though on the artist's own retina the head of a
person standing close may be much larger than the head of one standing deeper in the
group.  The size-constancy mechanism not only enlarges distant objects but also reduces
the apparent size of those that are too close~--- for instance a palm held up to the eye.

Experiments have shown that the property of size-constancy is nearly absolute, so
accurately does human consciousness reconstruct the true sizes of objects from the
retinal image given information about their distances; yet two factors were found
invariably to reduce or completely eliminate the size-constancy mechanism:
(1)~constancy was disrupted for distant objects;
(2)~constancy was not maintained even for nearby objects if those objects had few or
no depth cues.\footnote{%
  The absence of depth cues was established by the following method, for example:
  the apparent size of a luminous disk at various distances from the observer was
  assessed in absolute darkness, under monocular observation.}

Both these experimental facts accord fully with the theoretical picture sketched
above~--- information about distant objects has secondary biological value, and the
formation of three-dimensional perceptual space is possible only when depth cues are
used that are not contained directly in the retinal image.

The form-constancy mechanism is analogous in nature but concerns the \emph{form}
of objects.  Suppose a person looks at objects of simple shape~--- a circle, a square, and
the like (for instance, the surfaces of a round or square table).  If viewed at some
arbitrary angle, a circle will be perceived as an ellipse, a square as something
resembling a rhombus, and so on.  The action of the form-constancy mechanism
consists in this: if the person already knows the true form of the objects from prior
experience, round objects appear to him less oval than their retinal images; the angles
of the rhombs appear closer to~90° than in the images on the retina (i.e.\ closer to a
square), and so on.  Human consciousness endeavours to compensate not only for
distortions of relative size, but also for the distorted \emph{form} of objects arising when
they are projected onto the retina through the optical system of the eye.  It is well
known, for example, that a television screen can be viewed from the side at a
considerable deviation from the perpendicular to the screen plane, yet the apparent form
of the screen is preserved as the familiar horizontally elongated rectangle~--- even
though its retinal image may already have become square or even vertically elongated.

\srcpage{48}

The form-constancy mechanism can be studied quantitatively with comparative ease
by conducting experiments with objects of simple form (circle, square, etc.); this does
not mean, however, that it does not act on the retinal images of more complex objects.

Concerning the effectiveness of the form-constancy mechanism, the following should
be noted:
(1)~as already mentioned, its action is stronger the better the form of the object
contemplated is known (in advance, from prior experience);
(2)~it acts only at comparatively small distances and is completely ineffective for
objects sufficiently remote from the viewer.

The reasonableness of both conditions from a biological standpoint is sufficiently
evident.

Returning to the question of the two-stage character of human visual perception (to
the extent of detail needed for what follows), one may assert that the first stage is the
formation on the retina of an image of the external world by the optical system of the
eye, while the second is the reconstruction on that basis of three-dimensional perceptual
space, with subconscious use of depth cues and tactile-kinesthetic representations, by
means of the transformation of the retinal image by the size-constancy and
form-constancy mechanisms.  The latter, as it were, ``stretch'' and ``compress''
individual elements of the retinal image, subjecting to the most significant
transformations that portion of the retinal image corresponding to the region of the
person's immediate surroundings, which is of primary biological importance.

What has been said here may create an excessively simplified impression of how the
brain works during visual perception.  To clarify somewhat the character of the brain's
work (again only to the extent required for what follows) and to illustrate some of the
assertions made above, let us turn to the following example.

\figblock{14}{Axonometric images of a cube.  Centre: the wire-frame cube
  (14A)~--- no depth cues are effective and the spatial interpretation is ambiguous.
  Left and right: when the faces are opaque (14B, 14C), the occlusion cue
  resolves the ambiguity, giving the two possible solid orientations of the cube.}

In the central part of Fig.\,14 a cube is depicted using \emph{axonometric projection},
in which the property of parallelism is preserved (edges that are parallel in the real
cube are drawn as mutually parallel lines).  This drawing is remarkable in that only the
edges of the cube are shown (a cube made, for example, as a wire frame), so that the
depth cues listed above are ineffective.  Indeed, the binocular cues cannot assist correct
perception since a real cube is not presented to the eyes, but only its image.
Furthermore, the binocular depth cues when contemplating any picture of a
three-dimensional object will always be only a hindrance, since they will emphasize that
all points of the image are equidistant from the viewer~--- at the distance from the eyes
to the paper, canvas, or panel~--- and will thereby undermine the illusion of depicted
space (if the artist has set himself the task of conveying spatial depth).

\srcpage{49}

As for the monocular depth cues, an artist can only rely on those monocular cues listed
above.  Inspecting those five cues, it is easy to see that in this particular case none of
them can give the brain the necessary information; the only cue that might prove
useful~--- more distant objects or their parts are seen closer to the horizon~--- is not
absolute, for a part of an image being close to the horizon line may result not only from
its distance from the observer but also from an actual displacement in objective space of
the object or its part in the vertical direction.

This is illustrated by the other drawings of the same cube (14B, 14C), which show two
possible orientations of the cube (14A)~--- interpretations that became definite as soon
as the faces of the cube ceased to be transparent and the absolute depth cue of occlusion
came into play (nearby objects or their parts block more distant ones from view).

It should be said that occlusion is almost the only depth cue when axonometry is used as
a method of depicting three-dimensional objects.  Indeed, the other cue~--- diminution
of size with depth~--- is inapplicable in axonometry, since axonometry by definition
does not alter sizes with displacement in depth (the front and back edges of the cube
have the same size in the drawing); common sense suggests (and the next chapter will
demonstrate rigorously) that axonometry as a depiction method applies only to shallow
spaces~--- so that effects of the aerial-perspective type cannot be used.  As for the
casting of shadows, this depth cue is appropriate in axonometric depictions but plays a
subordinate role (one must first obtain depictions of the type 14B or 14C, and only then
cast shadows; casting shadows on the transparent cube 14A is pointless).

In the examples discussed, the talk was of the subconscious engagement of some depth
cue in the process of perception.  However, on more careful analysis one can also see
here the role of conscious awareness.  In the discussion of cube perception above it was
stated that what was depicted \emph{was} a cube, i.e.\ a certain definiteness was
introduced (through consciousness rather than the subconscious), and it remained only
to clarify which of the two faces drawn as squares was the front one.  This is precisely
what happened subconsciously in a simple glance at 14B and 14C\@.  Had the drawings
been identified not as a cube but as an empty box turned toward the viewer with the
open side, then, by interpreting this box as an empty room, one could say that 14B
depicts the upper right, and 14C the lower left, corner of the room.  Consequently,
visual perception is governed not only by the subconscious but also by conscious
awareness~--- and further: by looking at 14A and deliberately suggesting to oneself
what object is depicted there, one can easily make oneself see a flat pattern (for
instance a tile design for a floor), a cube, or an empty box, each in two positions,
without referring to the lower drawings at all.

\figblock{15}{A room and a die.  The two drawings have the same geometry (of
  the type 14C) but are immediately perceived as different three-dimensional objects,
  demonstrating the decisive role of prior experience in interpreting the retinal image.}

\srcpage{50}

It is therefore not surprising that in the study of the psychology of visual perception,
the proposition was advanced that when ``decoding'' some retinal image our
consciousness runs through all possible interpretations of it and, as a rule, selects the
one most frequently encountered in prior life experience, or instilled by training or
suggestion.  This fact has long been known: Helmholtz wrote roughly a hundred years
ago that ``the facts~\ldots\ show the profound influence of experience, training, and
habit on our perception.''

For the selection of one from several possible interpretations of the retinal image~---
and hence for the refinement of the geometric (spatial) characteristics of the visible
object~--- useful not only are the depth cues but also those cues that allow recognition
of the object.  Illustration 15 shows two diagrams of the type 14C; yet it is immediately
clear that the spatial images are different despite the geometric identity of their
structure.  Prior experience plays its role here, of course; a being from another planet
could not unambiguously identify these two images, since he would not know our houses
or our games.

Thus in reconstructing three-dimensional perceptual space from the retinal image, the
following are important:
(1)~depth cues;
(2)~a representation of the nature of the depicted object (its recognition);
(3)~the ``mental set'' of the observer, conveyed by training, suggestion, habit, and
personal interest.

In real life a combination of different depth cues acts (their relative importance varying
with the specific situation), together with an ensemble of recognizable objects that
mutually complement and refine the overall picture, and, of course, the observer's
mental set~--- a complex combination of representations drawn from life experience and
suggestion.  It must be said that the result of the combined action of all these factors is
so effective that people in their everyday lives almost never make mistakes; ``visual
illusion'' is the rarest exception in the practical activity of the human being.  The
situation becomes notably more complex when a person contemplates not an object but
its depiction in a picture, since in that case a number of cues contradict the depicted
subject (as already noted, the binocular depth cues), and therefore a relatively more
important role begins to be played by the group of factors listed in the last point~---
training, suggestion, personal interest.

\figblock{16}{An elongated parallelepiped (beam, long building, etc.) in axonometric
  projection.  Although the rear square face and the front square face are equal in the
  image, the rear one appears larger~--- an instance of the size-constancy mechanism
  acting on a picture.}

\srcpage{51}

Let us return once more to the size-constancy and form-constancy mechanisms.  Their
biological significance in practical life is evident and need not be discussed.  For the
artist, however, a question of primary importance is different: do these constancy
mechanisms operate when viewing pictures?  In pictures, after all, many depth cues not
only fail to assist but even impede the desired perception of the picture.

The answer to this question must be taken as affirmative.  Both mechanisms operate,
albeit in a weakened form, also in relation to depictions of objects.  Especially important
is the size-constancy mechanism, since it affects not a single object but the entire
structure of the depicted space.  To illustrate its action, Fig.\,16 shows the axonometric
image of an elongated parallelepiped (a beam, a long building, etc.).  The rear square
face appears larger than the front one, although in the image they are equal.  Here we
see a ``deception'' in the system by which the viewer perceives the image.  Since one of
the depth cues (the nearer blocks the farther) indicates a considerable remoteness of the
rear face, the brain attempts subconsciously to compensate for the corresponding
reduction of its size on the retina.  In axonometry no reduction of size occurs, and
therefore such compensation is unnecessary; yet the size-constancy mechanism
somewhat enlarges the perceived magnitude of the rear face's edges, and the
parallelepiped appears to widen with depth.

From what has been said, two conclusions may be drawn:
(1)~in some cases the constancy mechanisms can distort the picture in an undesired
direction;
(2)~with a judicious allowance for the effects associated with the action of the
constancy mechanisms, perception of a picture may be improved.

The second conclusion does not follow from Fig.\,16 alone, but is evident in itself.
